\bookchapter{Capítulo 33: Perspectivas futuras — a próxima etapa do celular}

\chapterintro{O celular não desaparecerá da noite para o dia, mas sua “posição de monopólio” pode ser dividida com uma nova geração de dispositivos. Neste capítulo, falamos de três coisas: como serão os dispositivos da próxima geração, como o celular se tornará “invisível” como base de uma nova identidade e as questões éticas que precisamos começar a pensar desde já.}

\booksection{AR/VR e computação espacial: o “segundo desktop” além do celular}

A competição da próxima geração de dispositivos já começou. Os que mais provavelmente dividirão o tempo de uso do celular são os dispositivos de AR/VR e computação espacial:

\begin{itemize}
  \item \textbf{Sucessor do Vision Pro}: O lançamento do Apple Vision Pro marcou a chegada oficial do conceito de "computação espacial", mas o preço (a partir de 29.999 RMB) e o peso da primeira geração ainda são barreiras. Espera-se que entre 2026 e 2027 seja lançado um Vision mais leve e acessível, levando a computação espacial de "brinquedo de desenvolvedor" para o "público inicial". Nesse cenário, o iPhone pode se tornar o parceiro de processamento e hub de dados dos óculos Vision.
  \item \textbf{Óculos AR}: Abrem um "segundo desktop" diante dos seus olhos, tornando a informação digital tão onipresente quanto o ar. Os óculos inteligentes da Meta em parceria com a Ray-Ban já venderam milhões de unidades. Na China, marcas como Rokid, XREAL e INMO também estão lançando óculos AR de consumo. Os óculos AR de 2025-2026 já conseguem fazer tradução em tempo real, HUD de navegação e assistente de IA por voz, mas ainda são limitados pela tecnologia de guia de onda óptica e duração da bateria — ainda faltam de 3 a 5 anos para "substituir o celular".
  \item \textbf{Anéis inteligentes e wearables}: Anéis inteligentes como Samsung Galaxy Ring e Oura Ring estão se popularizando como "sensores corporais" complementares do celular para monitoramento de saúde. O papel do celular está mudando de "terminal universal" para "hub central + múltiplos acessórios leves".
\end{itemize}

\booksection{Interface cérebro-computador: o pensamento como meio de interação}

O próprio pensamento se tornará um meio de interação, mas este caminho é mais longo do que se imagina:

\begin{itemize}
  \item \textbf{Últimos avanços do Neuralink}: O Neuralink de Elon Musk concluiu o primeiro implante humano em 2024. O primeiro paciente, Noland Arbaugh, conseguiu controlar o cursor do computador com a mente para jogar xadrez online. Atualmente, as BCIs (interfaces cérebro-computador) implantáveis são voltadas principalmente para assistência a pacientes com paralisia, com o objetivo de ajudá-los a "operar" novamente o mundo digital.
  \item \textbf{BCI não invasiva}: Várias empresas estão desenvolvendo tiaras não invasivas (como Kernel, OpenBCI) que leem sinais cerebrais sem cirurgia, mas com precisão muito inferior às implantáveis. Ainda há uma distância considerável para "digitar com a mente".
  \item \textbf{Perspectivas para o consumidor}: A interface cérebro-computador ainda está a 10-20 anos de distância da interação com celulares de consumo, mas a curto prazo pode ser aplicada em assistência a pessoas com deficiência e monitoramento de foco. Questões éticas e de privacidade ("a quem pertencem minhas ondas cerebrais?") se tornarão foco de debate antes mesmo do amadurecimento da tecnologia.
\end{itemize}

\booksection{A “invisibilidade” do celular e a nova base de identidade}

Outra direção futura de evolução do celular talvez seja a “invisibilidade” — transformar-se em um “agente oculto” onipresente, tornando-se a base única da sua identidade digital:

\begin{itemize}
  \item \textbf{Autenticação de dois fatores na nuvem}: Concluir instantaneamente a autenticação de "você é você", acessando qualquer serviço sem precisar digitar códigos de verificação manualmente.
  \item \textbf{Herança digital como configuração de fábrica}: Ao ativar um novo dispositivo, predefinir um plano de sucessão de ativos digitais, tornando o "pós-vida" uma configuração padrão de fábrica.
  \item \textbf{De "centrado no dispositivo" para "centrado na identidade"}: Seu cofre de dados pessoais na nuvem é o centro; o hardware é apenas o ponto de acesso. Ao trocar de dispositivo, sua identidade e dados sempre o acompanham.
\end{itemize}

\booksection{Reflexão ética: ainda somos donos do celular?}

Uma pergunta fundamental sempre existe: somos “deuses” empoderados pela tecnologia ou “cérebros em cubas” alimentados por algoritmos?

\begin{itemize}
  \item \textbf{Ética da IA}: Quando a IA no dispositivo te conhece melhor do que você mesmo, suas "escolhas" são vontade autônoma ou resultado de direcionamento? Quando um agente de IA faz pedidos e responde mensagens por você, onde estão a responsabilidade e os limites?
  \item \textbf{Soberania de dados}: Seus dados de saúde, dados cerebrais e dados biométricos pertencem a você, ao fabricante do dispositivo ou ao provedor de nuvem? A Lei de Proteção de Informações Pessoais e o GDPR da UE são apenas o ponto de partida — a implementação real do "direito à portabilidade de dados" e do "direito ao esquecimento" ainda tem um longo caminho pela frente.
  \item \textbf{Exclusão digital}: Celulares com IA, telas dobráveis e comunicação via satélite estão ampliando a distância entre quem "pode pagar" e quem "não pode pagar". O acesso equitativo à infraestrutura digital se tornará uma nova questão social — tão crucial quanto "saber ou não usar a internet" era há 20 anos.
  \item Quando o algoritmo sabe melhor do que você o que você quer ver, quanto do seu "gostar" é espontâneo?
  \item Quando você acende a tela inconscientemente a cada poucos minutos, isso é seu hábito ou um "comportamento viciante" projetado?
\end{itemize}

\booksection{Mantenha-se conectado, mas também independente — sua filosofia de vida digital}

Independentemente de como a tecnologia evolua, os seguintes princípios de filosofia de vida digital continuarão válidos daqui a 10 anos:

\begin{enumerate}
  \item Seja o distribuidor de recursos, não o receptor passivo de conteúdo empurrado.
  \item Deixe a automação ser seu fiel "copiloto".
  \item Construa sua "retaguarda digital" sólida.
  \item Realize periodicamente o "desapego digital" e o "jejum de redes sociais".
  \item Coloque sempre a segurança e a privacidade como pré-requisitos de toda conveniência.
\end{enumerate}

\begin{tcbpractice}{Experimente: Elabore seus "Termos de Serviço de Uso de Dispositivos Digitais Pessoais"}
Escreva um simples “termo de serviço” para você e seu celular (e futuros óculos AR, assistente de IA) para definir o relacionamento entre vocês. Por exemplo:

\begin{itemize}
  \item Artigo 1: Meu tempo de atenção me pertence. Você deve entrar ativamente no estado "Não Perturbe" quando eu estiver trabalhando ou com a família.
  \item Artigo 2: Meus dados pessoais são minha propriedade. Sua principal responsabilidade é a criptografia segura e a sincronização conveniente.
  \item Artigo 3: Nosso objetivo comum é aumentar a eficiência e a experiência de vida. Se você me induzir ao vício por meio de design persuasivo, tenho o direito de desinstalá-lo.
\end{itemize}

\end{tcbpractice}

\begin{tcbsummary}{Resumo do capítulo}
\begin{itemize}
  \item Vision Pro e óculos AR são novas portas de computação além do celular, mas dificilmente substituirão o celular a curto prazo
  \item A interface cérebro-computador está a 10-20 anos do consumidor; questões éticas precedem o amadurecimento técnico
  \item O celular está evoluindo para a "invisibilidade", do centrado no dispositivo para o centrado na identidade
  \item Ética da IA, soberania de dados e exclusão digital são temas centrais para os próximos 10 anos
  \item Manter-se conectado e também independente é a filosofia imutável da vida digital
\end{itemize}

\end{tcbsummary}
